%==================================================================
%  Quarto -> Beamer 用プリアンブル（docs/SlideRules.md 準拠）
%  econ.lua がこのファイルをヘッダに入れる。その直前に、YAML の
%  切り替えオプションから作った設定マクロが入る:
%    \econNoSectionOutline … section-outline: false のとき定義される
%    \econBibFile          … 文献ファイル（bib-file:、既定 ref.bib）
%  テーマの起動（\usetheme）は Quarto(YAML) 側ではなく、ここで完全に制御する。
%==================================================================

%========================
%  パッケージ
%========================
\usepackage{xcolor}
\usepackage{colortbl}      % \cellcolor（表セルの塗り）用
\usepackage{graphicx}
\usepackage{tabularx}
\usepackage{multirow}
\usepackage{booktabs}
\usepackage{amsmath, amssymb}
\usepackage{pgfplots}
\pgfplotsset{compat=1.18}
\usepackage{tikz}
\usetikzlibrary{shapes.geometric, arrows.meta, calc, positioning}
\usepackage{tcolorbox}
\tcbuselibrary{skins}
\usepackage{caption}
\captionsetup[figure]{labelformat=empty, skip=2pt, font=small}
\usepackage{appendixnumberbeamer}

%========================
%  色の定義
%  役割: blue=構造＋文中強調（単独可） / red=文中強調・対比の一方 / orange=box・表の強調（\ybox 等）
%  文中強調は \textcolor{mainred}{...} / \textcolor{mainblue}{...} を直接使う
%  （\alert はredのみでblueに使えないため、red/blue共通の書き方として統一）。
%  注意: orange はテキスト色として単体で使うと白背景でのコントラストが
%  弱く（実測コントラスト比 約1.6:1）、プロジェクター投影で特に読みにくい。
%  そのため文中強調には使わず、枠線・薄い塗りの box/表セル用途に限定する。
%========================
\definecolor{textblack}{RGB}{42,46,51}
\definecolor{textgray}{RGB}{141,142,143}
\definecolor{bgwhite}{RGB}{255,255,255}
\definecolor{mainblue}{RGB}{54,136,214}     % 構造の主役（frametitle帯・箇条書き点など）
\definecolor{subblue}{RGB}{227,241,255}     % frametitle帯・定義ボックス地（最も淡いブルー）
\definecolor{mainred}{RGB}{210,48,55}       % 文中強調・対比の一方（白背景コントラスト比 約5.0:1）
\definecolor{subred}{RGB}{249,183,185}
\definecolor{mainorange}{RGB}{240,199,74}   % box・表セルの強調（\ybox・cellcolor 用。文字色には使わない）
% 中間トーン（ナビバー帯・block body など）は主色から自動計算で導出
\colorlet{midblue}{mainblue!35!white}

%========================
%  テーマ（metropolis + 上部ナビバー miniframes）
%  ※ Quarto YAML では theme を指定せず、ここで完全制御する
%========================
\usetheme{metropolis}

\setbeamercolor{normal text}{fg=textblack, bg=bgwhite}
\setbeamercolor{background canvas}{bg=bgwhite}
\setbeamercolor{title}{fg=textblack}
\setbeamercolor{subtitle}{fg=mainblue}
\setbeamercolor{institute}{fg=mainred}
\setbeamercolor{frametitle}{fg=textblack, bg=subblue}
\setbeamercolor{block title}{bg=mainblue, fg=white}
\setbeamercolor{block body}{bg=midblue}
\setbeamercolor{alerted text}{fg=mainred}
\setbeamercolor{block title alerted}{bg=mainred, fg=white}
\setbeamercolor{block body alerted}{bg=subred}
\setbeamercolor{block title example}{bg=mainblue, fg=white}
\setbeamercolor{block body example}{bg=midblue}
\setbeamercolor{progress bar}{fg=mainred, bg=subred}
\setbeamercolor{palette primary}{bg=mainblue, fg=white}
\setbeamercolor{palette secondary}{bg=midblue, fg=textblack}
\setbeamercolor{palette tertiary}{bg=mainred, fg=white}
\setbeamercolor{palette quaternary}{bg=subred, fg=textblack}

% 上部セクション進行バー（現在=黒/太字、他=グレー、丸印で位置表示）
\useoutertheme[subsection=false]{miniframes}
\setbeamerfont{section in head/foot}{series=\bfseries}
\setbeamercolor{section in head/foot}{fg=textblack, bg=midblue}
\addtobeamertemplate{frametitle}{\vskip -0.01cm}{}
\setbeamertemplate{navigation symbols}{}

% Outline スライド（\tableofcontents[currentsection]）：現在のセクションは太字・blue、
% それ以外は既定のシェード（薄いblue、グレーではない）のまま変更しない
\setbeamerfont{section in toc}{series=\bfseries}
\setbeamercolor{section in toc}{fg=mainblue}

% タイトル下の区切り線の色（metropolis のバージョンにより反映されない場合あり。
% 反映されないときは既定のグレーのまま＝十分自然なのでそのままでよい）
\setbeamercolor{title separator}{fg=mainblue}

% 発表者オリジナルの意見・構成のスライドを示す強調色。
% 注意: \setbeamercolor{frametitle}{...} を ## 見出しの前後の生LaTeXブロックで
% 切り替える方式は、pandoc がそのブロックを「前フレームの本文末尾」に
% 配置し、かつ beamer の frame 環境が内部でグループ化されているために
% \end{frame} で変更が巻き戻ってしまい機能しない（要注意点）。
% そのため、対象スライドは見出し自体を \textcolor{mainred}{...} で
% 直接色付けする方式を採用している（template.qmd の該当 ## 見出しを参照）。

% \appendix が part page（"Appendix" ボックス）を挿入しないよう無効化
\AtBeginPart{}

% 箇条書きの記号色を blue に（L1 の ●）
\setbeamercolor{itemize item}{fg=mainblue}
\setbeamercolor{itemize subitem}{fg=mainblue}
\setbeamercolor{itemize subsubitem}{fg=mainblue}
\setbeamercolor{enumerate item}{fg=mainblue}

% 各セクションの冒頭に、現在位置を強調した目次（Outline）を自動挿入（非該当セクションは淡色）。
% YAML で section-outline: false のときは入れない（短い発表用）。
% このとき \AtBeginSection を空にするのは、metropolis 既定のセクション区切りページ
% （これも \AtBeginSection で入る）まで出てこないようにするため。
% 注意: tableofcontents の各項目は内部的に \hyperlink なので、
% \hypersetup の linkcolor（=textblack）の方が section in toc の色設定より
% 優先されてしまい、何もしないと常に黒字になる。ここでローカルに
% linkcolor を上書きして、狙い通り現在セクションが blue になるようにする。
\ifdefined\econNoSectionOutline
  \AtBeginSection[]{}
\else
  \AtBeginSection[]{%
    \begin{frame}{Outline}
      \begingroup
      \hypersetup{linkcolor=mainblue}
      \tableofcontents[currentsection]
      \endgroup
    \end{frame}}
\fi

%========================
%  フォント
%  Latin=Noto Sans / 日本語=Noto Sans JP（欧文と同ファミリーで統一）
%  ※ Yu Gothic は通常ウェイトが細くプロジェクター投影で見づらいため、
%     欧文と揃うNoto Sans JPを使う。通常=Medium／太字=SemiBoldを明示指定しないと
%     Noto Sans JP Medium 単体では太字シェイプが無く \textbf が効かなくなる点に注意。
%========================
\usepackage[sfdefault]{noto}
\usepackage[match]{luatexja-fontspec}
\setmainjfont{Noto Sans JP}[UprightFont={* Medium},BoldFont={* SemiBold}]
\setsansjfont{Noto Sans JP}[UprightFont={* Medium},BoldFont={* SemiBold}]

%========================
%  tikz スタイル（決定木・DAG 用）
%========================
\tikzset{
  dag_node/.style={
    draw=textblack, ellipse,
    minimum width=2.0cm, minimum height=1.0cm,
    thick, align=center, fill=bgwhite, font=\small
  },
  box_node/.style={          % 決定木の通常ノード（blue）
    draw=mainblue, fill=midblue, rounded corners=2pt,
    minimum width=2.4cm, minimum height=0.9cm, align=center, font=\small
  },
  key_node/.style={          % 強調したい結論ノード（red系で差別化）
    draw=mainred, fill=subred, rounded corners=2pt,
    minimum width=2.4cm, minimum height=0.9cm, align=center, font=\small\bfseries
  },
  arrow/.style={
    ->, >=Stealth, thick, line width=1.0pt,
    shorten >=2pt, shorten <=2pt
  }
}

%========================
%  カスタム部品（SlideRules.md の型に対応）
%========================
% (1) 数式の変数定義ゾーン（淡いblue地のボックス）
\newtcolorbox{defbox}{
  colback=subblue, colframe=subblue,
  boxrule=0pt, arc=3pt, left=8pt, right=8pt, top=6pt, bottom=6pt,
  before skip=6pt, after skip=6pt
}

% (2) 表・数値の「オレンジ枠」強調。
%     枠線・薄い塗りとしての使用なので、文字色ほどコントラストは問題にならない。
%     セル背景で強調したい場合は表内で \cellcolor{mainorange!25}{...} を使う。
%     注意: \ybox は "on line"（改行しない箱）なので短い語句専用。
%     文章まるごと強調したいときは下の keytakeaway 環境（複数行対応）を使う。
\newtcbox{\ybox}{on line, colframe=mainorange, colback=mainorange!25,
  boxrule=1pt, arc=2pt, boxsep=1.5pt, left=2pt, right=2pt, top=1pt, bottom=1pt}

% (2b) 文章単位の「まとめ・要点」強調（折り返し対応、オレンジ地）。
%      表・数値の強調に加え、議論の転換点・核心の一文を格上げする用途にも使ってよい
%      （多用しない＝1デッキで数回程度、本当に格上げしたい一文だけに絞る）。
\newtcolorbox{keytakeaway}{
  colback=mainorange!25, colframe=mainorange,
  boxrule=1pt, arc=2pt, left=6pt, right=6pt, top=4pt, bottom=4pt,
  before skip=4pt, after skip=4pt, fontupper=\small
}

% (3) Table / Figure ラベル（blue 太字）
\newcommand{\tlabel}[1]{\textcolor{mainblue}{\textbf{#1}}}

% (3b) 図表下の注釈・箇条書き項目の出典（グレー・小さめ）。行頭の生波括弧を避けるためコマンド化。
\newcommand{\slidenote}[1]{{\footnotesize\color{textgray}#1}}

% (3c) \slidenote よりさらに小さい注（表の下のSE・有意水準の注記など、詰めて書きたい箇所用）。
\newcommand{\tinynote}[1]{{\tiny\color{textgray}#1}}

% (3d) 箇条書き項目の出典。(Author, Year) が途中で折り返さないよう \mbox で包む。
%      使い方: 項目本文の末尾に \srcite{key}
\newcommand{\srcite}[1]{\slidenote{\mbox{\parencite{#1}}}}

% (4) 図のプレースホルダ（実画像が未配置でもレンダリング可能にする）
%     本番では \includegraphics{ファイル名} に置き換える。
\newcommand{\figplaceholder}[2][7cm]{%
  \begin{tikzpicture}
    \node[draw=textgray, dashed, rounded corners=4pt, fill=black!2,
          minimum width=#1, minimum height=3.4cm,
          align=center, text=textgray, font=\small]
      {\itshape 図をここに配置\\[2pt] \ttfamily #2};
  \end{tikzpicture}}

%========================
%  参考文献（biblatex / APA + biber）
%  本文中は \textcite{key} / \parencite{key}、文献ページは \printbibliography を使う。
%  ※ citeproc: false（_extension.yml）で pandoc の引用処理を止め、biblatex に一本化。
%  ※ 文献ファイルは YAML の bib-file:（既定 ref.bib）。論文と同じ .bib を共有したいとき
%     bib-file: ../paper/refs.bib のように指定する。
%========================
\usepackage{csquotes}
\usepackage[style=apa, backend=biber]{biblatex}
\DeclareLanguageMapping{english}{english-apa}
\providecommand{\econBibFile}{ref.bib}
\addbibresource{\econBibFile}
\AtEveryBibitem{
  \clearfield{doi}\clearfield{url}\clearfield{isbn}\clearfield{issn}
  \clearfield{note}\clearfield{month}\clearfield{issue}\clearfield{number}
  \clearfield{pmid}\clearfield{pmcid}\clearfield{abstract}
}
% 文献リストの見た目：先頭アイコンを消し、タイトルを通常ウェイトに
\setbeamertemplate{bibliography item}{}
\setbeamerfont{bibliography entry title}{series=\mdseries}

%========================
%  ハイパーリンク
%========================
\hypersetup{colorlinks=true, urlcolor=mainblue, linkcolor=textblack, citecolor=textblack}
